\begin{textsample}{POS Dim 1 – vem\_ed\_22 – Score 2.00 – vem\_ed\_22/t112.txt}  \label{ex:f1_pos_177}
Equipe dos PCIs na DUOC UC (Chile) Osvaldo Succi Junior, coordenador dos Projetos Colaborativos Internacionais da Unidade do Ensino Superior de Graduação (PCIs / Cesu) do Centro Paula Souza, e Regiane Souza Camargo Moreira, responsável pelos PCIs em espanhol, \textbf{foram} convidados pela Vice-Reitoria Acadêmica do Instituto Profesional DUOC UC (Chile) para oferecer workshops sobre a metodologia COIL (Collaborative Online International Learning) de Intercâmbios Virtuais e participar de seminários da “1ª Semana de Internacionalização em Casa”.
O evento ocorreu na instituição chilena entre 27 de novembro e 1º de dezembro de 2023, no centro de educação continuada da DUOC UC em Santiago e no campus Valparaíso.
As atividades \textbf{foram} organizadas para diversos perfis de público acadêmico: workshops para professores iniciantes em COIL (72 participantes), experientes na metodologia (12 participantes), masterclass sobre Internacionalização em Casa e seminário para grupo administrativo sobre essa temática.
O primeiro Intercâmbio Virtual (ou PCI, como \textbf{é} chamado no Centro Paula Souza) \textbf{foi} realizado na Fatec Americana, em 2013, precedendo \textbf{todas} as instituições de ensino \textbf{superior} brasileiras.
Em uma década, os mais de 500 PCIs envolveram 50 Fatecs, 9.934 alunos e 49 instituições de ensino \textbf{superior} em 18 países.
Hoje, o Centro Paula Souza está entre as principais instituições do mundo em Intercâmbios Virtuais do tipo COIL.

% matched lemmas: ser, superior, todo
\end{textsample}
